% Seminar 2, page 1 - manuscript dated 10-Oct-1994.
% The final digit of the year is partly hidden by the damaged corner;
% 1994 is inferred from the chronology of the notes.
% Fragment for \input: no preamble and no forced page break.
\subsection*{Seminar 2 - Page 1}

\begin{center}
\underline{\textbf{ALGEBRA}}\\[2pt]
\textit{(Seminar 2)}
\end{center}
\begin{flushright}
\textit{10-Oct-1994}
\end{flushright}

\noindent Monday:\quad algebra seminar, 3 hours;\quad geometry seminar, 3 hours.

\bigskip
\noindent\fbox{1}\quad Let $n\in\mathbb{N}^{*}$. Compute
\[
 d(n)=\left|\left\{
 f:\{1,2,\ldots,n\}\longrightarrow\{e_1,e_2,-e_1,-e_2\}
 \ \middle|\ \sum_{k=1}^{n}f(k)=0
 \right\}\right|,
\]
\noindent where $e_1=(1,0)$, $e_2=(0,1)$.

\medskip
\noindent\underline{Solution:}
\[
 d(n)=\begin{cases}
 0, & n\text{ odd},\\[2pt]
 \text{to be calculated}, & n\text{ even}.
 \end{cases}
\]

\noindent Let $n=2m$. Then
\[
 d(2m)=\sum_{j=0}^{m}
 \binom{2m}{j}\binom{2m-j}{j}\binom{2m-2j}{m-j}.
\]
\[
\begin{aligned}
 d(2m)
 &=\sum_{j=0}^{m}
 \frac{(2m)!}{j!(2m-j)!}\,
 \frac{(2m-j)!}{j!(2m-2j)!}\,
 \frac{(2m-2j)!}{(m-j)!(m-j)!}\\[4pt]
 &=\sum_{j=0}^{m}\frac{(2m)!}{(m!)^2}
 \left(\frac{m!}{j!(m-j)!}\right)^2\\[4pt]
 &=\frac{(2m)!}{(m!)^2}\sum_{j=0}^{m}\binom{m}{j}^{2}\\[4pt]
 &=\frac{(2m)!}{(m!)^2}\binom{2m}{m}\\[4pt]
 &=\frac{(2m)!}{(m!)^2}\cdot\frac{(2m)!}{m!\,m!}
 =\frac{(n!)^2}{\left[\left(\frac n2\right)!\right]^4}.
\end{aligned}
\]
% The manuscript has a small counting sketch beside the first sum,
% labelled j and 2m-2j; the corresponding combinatorial formula is retained.

\bigskip
\noindent\fbox{2}\quad\underline{(Characterisation of injective functions)}

\noindent Let $f:A\longrightarrow B$. Show that the following are equivalent:
\begin{enumerate}
\item[a)] $f$ is injective.
\item[b)] $f$ has a retraction, i.e. there exists $r:B\longrightarrow A$
 such that $r\circ f=1_A$.
\item[c)] $f$ is a monomorphism: for every set $X$ and
 $f_1,f_2:X\longrightarrow A$,
 \[
 f\circ f_1=f\circ f_2\ \Longrightarrow\ f_1=f_2.
 \]
 \[
 \begin{tikzcd}[column sep=large]
 X \arrow[r, shift left=.55ex, "f_1"]
   \arrow[r, shift right=.55ex, "f_2"']
 & A \arrow[r, "f"] & B
 \end{tikzcd}
 \]
\end{enumerate}
% A marginal arrow links (a) and (c) graphically, not as a separate proof.
% The equivalence (a) <=> (b) presupposes A is nonempty (or a suitable
% convention for empty sets); the manuscript does not state this condition.
